% RTX Software Engineering Co-Op Cover Letter -- Wilsonville, OR
\documentclass[letterpaper,11pt]{letter}

\usepackage[empty]{fullpage}
\usepackage[hidelinks]{hyperref}
\usepackage[english]{babel}
\usepackage{iftex}
\ifPDFTeX
  \input{glyphtounicode}
  \pdfgentounicode=1
\fi

\addtolength{\oddsidemargin}{-0.5in}
\addtolength{\evensidemargin}{-0.5in}
\addtolength{\textwidth}{1in}
\addtolength{\topmargin}{-0.5in}
\addtolength{\textheight}{1.0in}
\raggedright
\setlength{\parindent}{0pt}
\setlength{\parskip}{10pt}

\begin{document}

\begin{flushleft}
\textbf{\large Prabhas Gade}\\
832-759-3256 $|$ \href{mailto:prabhasgade333@gmail.com}{prabhasgade333@gmail.com}\\
\href{https://linkedin.com/in/prabhasgade}{linkedin.com/in/prabhasgade} $|$
\href{https://github.com/PrabhasG550}{github.com/PrabhasG550} $|$
\href{https://studiorain.xyz}{studiorain.xyz}
\end{flushleft}

\vspace{6pt}
\begin{flushleft}
RTX Software Engineering Co-Op Recruiting\\
Wilsonville, OR -- Onsite Co-Op (May/June--December)
\end{flushleft}

Dear RTX Recruiting Team,

I am applying for the Software Engineering Co-Op in Wilsonville because I want to spend a full May/June--December session designing, debugging, integrating, and verifying software in an aerospace environment where correctness and safety matter. As a Computer Science (STEM) student at the University of Texas at Dallas with a 3.5 GPA and expected graduation in May 2028, I can remain enrolled through the co-op session and fully engage in the mentorship opportunity, including taking a semester away from coursework as required.

My strongest preparation for this role is industrial and embedded software development. At Vinson Process Controls, I designed and developed software from functional and interface requirements for an edge-computing RTU using React, Node.js, and MySQL in an Agile team. I performed system integration across UI, services, and persistence, and created requirement-based test procedures and automated unit tests before production deployment. Debugging hardware--software interaction issues in device environments taught me to be disciplined about reproduction, isolation, and verification---habits that transfer directly to laboratory and simulation work on flight-critical systems.

I also bring embedded and C\# experience that aligns with RTX's preferred stack. In MudraDetection, I developed on-device computer-vision software on an NVIDIA Jetson Orin Nano, optimized inference with ONNX/TensorRT, and verified performance through structured evaluation procedures and coverage of held-out and live test scenarios. In Project Pawrkour, I used C\# and Visual Studio to implement and debug real-time gameplay systems and camera/collision integration. Across full-stack work with React, Angular, Node.js, Django, Java/Eclipse, Python, and Git, I strengthened automated testing and clear communication of technical findings with engineers and stakeholders.

I am excited to learn Head-Up Display and camera systems, deepen my aerospace software practices, and contribute to software development, unit testing, test procedures, system integration, and formal verification support on your team. Thank you for your consideration. I would welcome the opportunity to bring careful engineering and a strong learning mindset to RTX.

Sincerely,\\
Prabhas Gade

\end{document}
